\movedhead{Is the compression regression to the mean?}
\label{sec:rtm}\label{app:rtm}
A falling spread has an innocent explanation that would account for the whole
result. Each family effect is estimated from twelve pairs with five trials an
arm, so each carries sampling error; sampling error inflates the spread of
the estimates above the spread of the truth, and it inflates an absolute
value most where the true effect is near zero. If that error were larger at
\texttt{low} effort than at \texttt{high}, the spread would fall with effort
even if nothing about the model's bias had changed. The analysis in
\texttt{analysis/robustness.py} tests this three ways. None of it is
pre-registered.

\paragraph{Removing the measurement error.} The variance of the nine
estimates is the variance of the truth plus the mean sampling variance of an
estimate, so the second can be subtracted. Two floors bracket the answer.
The narrow floor counts only the trial resampling variance of each pair's own
difference, which is unambiguously measurement error; subtracting it leaves
the spread almost unchanged and the fall survives in all six combinations,
for instance Haiku's continuous spread at \texttt{low} moving from 12.91 to
12.54 and its categorical spread from 0.215 to 0.206. The wide floor instead
counts the whole standard error of the family mean across its twelve pairs,
which also removes the genuine item to item variation within a family; it is
too harsh, because a cue whose effect differs across items is a fact about
the model rather than a measurement fault. Under that floor the fall survives
in three of six: Haiku on both measures and Sonnet on the continuous measure.
The other three cells have no between-family spread left to compress at any
effort level, which is itself worth stating plainly. In Opus and in Sonnet on
the categorical measure the nine families are indistinguishable from one
another once item variation is charged against the estimate, so for those
cells the design can establish that the raw spread falls but not that a real
between-family spread existed to begin with.

\paragraph{The noise runs against the finding in two models.} The trial-level
noise floor is not constant across effort, and its direction matters. In
Haiku it rises from 0.062 to 0.073 on the categorical measure and from 3.10
to 3.53 on the continuous one, and in Sonnet from 0.023 to 0.056, so in those
models the \texttt{high} effort estimate is the noisier one. Noise inflates a
spread, so this biases the comparison against the result, and the spread
falls anyway. In Opus the floor falls with effort, from 0.039 to 0.024 and
from 1.49 to 1.16, which is the direction that would help the result; Opus is
therefore the model whose compression rests most on the correction, and it is
also the one where the wide floor leaves nothing.

\paragraph{A placebo.} Shuffling the effort label among the responses to the
same item within a model destroys any real effect of effort and leaves
everything else intact. Over 2{,}000 shuffles the spread falls in all six
combinations 6.25\% of the time, so the \textit{count} of six out of six is
not on its own strong evidence and should not be reported as though it were.
The magnitudes are a different matter. Every one of the six observed changes
is more negative than almost every shuffle: $p < 0.0005$ for Haiku on both
measures, where no shuffle in 2{,}000 reached the observed change, $p =
0.0065$ and $p = 0.0010$ for Sonnet, and $p = 0.0080$ and $p = 0.0100$ for
Opus. The reason the count is weak while each magnitude is strong is that the
six cells are positively correlated under the placebo, so a shuffle that
makes one fall tends to make several fall.

\paragraph{Leaving a family out.} Dropping each family in turn and recounting
leaves six of six intact in eight of the nine cases and five of six when
framing is dropped, so the result is not carried by any single cue.

Taken together: the compression is real in Haiku, where it is large and
survives every correction; real in Sonnet on the continuous measure; and in
Opus and Sonnet's categorical measure it is a fall in a spread that was never
established to exceed measurement error. The claim in the title of this
section is answered no for Haiku and not fully excluded for Opus. The paper
states the mechanism as a finding about the weakest model and a conjecture
about the others, and the remedy is more items per family rather than more
trials per item, since the limiting quantity is the twelve pairs.
